% 2-8-longterm.tex
%  Budget: 2pp.
%  Rules: every numeral in \lr{}; cite only from references.bib;
%  em-dash must be \lr{---} (bare --- is a tofu box in B Nazanin).

\section{افق‌های بلندمدت و پیش‌بینی احتمالاتی}

دو امتداد از مسئله مرکزی این پایان‌نامه وجود دارد که آگاهانه خارج از
محدوده آن قرار گرفته‌اند، اما مرور آن‌ها لازم است: افق‌های بلندتر از
روز-پیش، و پیش‌بینی احتمالاتی به‌جای نقطه‌ای.

\subsection*{افق‌های میان‌مدت و بلندمدت}

یوسفی و همکاران \cite{yousefi_long-term_2019} پیش‌بینی بلندمدت قیمت برق را
با روش‌های یادگیری ماشین بررسی می‌کنند و انگیزه اقتصادی متفاوتی را برجسته
می‌سازند: پروژه‌های بزرگ انرژی، که سرمایه آن‌ها از ابتدا قفل می‌شود و
بازگشت آن دهه‌ها طول می‌کشد، به برآورد قیمت در افق‌های بسیار بلند نیاز
دارند. این با انگیزه پیش‌بینی روز-پیش \lr{---} تنظیم پیشنهاد قیمت برای فردا
\lr{---} تفاوت ماهوی دارد.

گلاسی و همکاران \cite{ghelasi_day-ahead_2025} پل میان این دو افق را
می‌سازند و مدل‌های اقتصادسنجی را از افق روز-پیش تا افق‌های میان‌مدت و
بلندمدت گسترش می‌دهند. آنان تصریح می‌کنند که مدل‌های مقاوم برای این
افق‌ها به دلیل پویایی‌های در حال تحول بازار و نوسان‌پذیری قیمت اهمیت
فزاینده‌ای یافته‌اند.

\subsection*{پیش‌بینی احتمالاتی}

زیل و همکاران \cite{ziel_probabilistic_2018} پیش‌بینی احتمالاتی میان‌مدت و
بلندمدت را بررسی می‌کنند. استدلال آنان روشن است: آزادسازی بازارها و توسعه
منابع تجدیدپذیر، عدم‌قطعیت فزاینده‌ای درباره حرکت آینده قیمت پدید آورده
است، و در حالی که حجم زیادی از مقالات به پیش‌بینی کوتاه‌مدت پرداخته‌اند،
فرصت‌های تازه‌ای برای بازیگران بازار در مدل‌سازی این عدم‌قطعیت وجود دارد.

تفاوت بنیادی این رویکرد با پیش‌بینی نقطه‌ای آن است که خروجی مدل، یک عدد
نیست بلکه یک توزیع یا مجموعه‌ای از چندک‌هاست. این خروجی برای مدیریت ریسک
مناسب‌تر است، اما معیارهای ارزیابی متفاوتی می‌طلبد و بنابراین مستقیماً با
اعداد پیش‌بینی نقطه‌ای قابل مقایسه نیست.

\subsection*{چرا این دو امتداد در پژوهش حاضر دنبال نشدند}

هر دو حوزه، امتدادهای مشروعی برای این کار هستند و در فصل پنجم به‌عنوان
پیشنهاد پژوهش آینده مطرح می‌شوند. دلیل کنارگذاشتن آن‌ها در پژوهش حاضر،
روش‌شناختی است و نه کم‌اهمیت بودن آن‌ها.

پژوهش حاضر بر یک مقایسه کنترل‌شده میان پنج خانواده مدل بنا شده است. افزودن
افق بلندمدت یا خروجی احتمالاتی، به معنای افزودن مجموعه‌ای دیگر از معیارهای
ارزیابی و مجموعه‌ای دیگر از تصمیم‌های طراحی است؛ و هر تصمیم طراحی اضافی،
درجه آزادی تازه‌ای می‌افزاید که اعتبار مقایسه اصلی را رقیق می‌کند. با
توجه به اینکه انگیزه محوری این پایان‌نامه، چنان‌که در بخش \lr{2-1} آمد،
دقیقاً پاسخ به ضعف \emph{اعتبار مقایسه} در این ادبیات است، گسترش دامنه با
همان انگیزه در تعارض قرار می‌گرفت.

با این حال، یک پیوند مفهومی میان بخش حاضر و یافته‌های این پژوهش وجود دارد.
تأکید زیل و همکاران بر عدم‌قطعیت، و تأکید گلاسی و همکاران بر پویایی‌های
در حال تحول بازار، هر دو به این نکته اشاره دارند که یک پیش‌بینی نقطه‌ای
بدون بیان دامنه اعتبار آن، اطلاعات ناقصی به تصمیم‌گیرنده می‌دهد. آزمون
برون‌توزیع این پژوهش، پاسخی از جنس دیگر به همین نگرانی است: به‌جای
کمّی‌کردن عدم‌قطعیت درون یک دوره، نشان می‌دهد که وقتی دوره تغییر می‌کند
چه بر سر پیش‌بینی می‌آید.
